% paper93-copilot-lifecycle.tex — The Copilot Lifecycle:
%   Connection, Health, and Trust in JuGeo's Kernel
% Paper 93 in the Judgment Geometry series.
% Compile: pdflatex paper93-copilot-lifecycle && pdflatex paper93-copilot-lifecycle

\documentclass[11pt]{article}
\usepackage{lmodern}
% jugeo-common.tex — shared preamble for all Judgment Geometry papers
% Include via: % jugeo-common.tex — shared preamble for all Judgment Geometry papers
% Include via: % jugeo-common.tex — shared preamble for all Judgment Geometry papers
% Include via: \input{jugeo-common}

% ── Packages ────────────────────────────────────────────────────────────
\usepackage{amsmath,amssymb,amsthm,mathtools}
\usepackage{tikz-cd}
\usepackage{listings}
\usepackage{xcolor}
\usepackage{xspace}
\usepackage{booktabs}
\usepackage{multirow}
\usepackage{enumitem}
\usepackage[expansion=false]{microtype}
\usepackage{hyperref}
\usepackage{cleveref}
% \usepackage{thmtools}  % disabled: not available in this TeX installation
\usepackage{float}
\usepackage{caption}
\usepackage{subcaption}
\usepackage{tabularx}
\usepackage{stmaryrd}       % \llbracket, \rrbracket
\usepackage{bbm}            % \mathbbm{1}

% ── Page geometry ───────────────────────────────────────────────────────
\usepackage[margin=1in]{geometry}

% ── Theorem environments ────────────────────────────────────────────────
\theoremstyle{plain}
\newtheorem{theorem}{Theorem}[section]
\newtheorem{lemma}[theorem]{Lemma}
\newtheorem{proposition}[theorem]{Proposition}
\newtheorem{corollary}[theorem]{Corollary}

\theoremstyle{definition}
\newtheorem{definition}[theorem]{Definition}
\newtheorem{example}[theorem]{Example}
\newtheorem{construction}[theorem]{Construction}

\theoremstyle{remark}
\newtheorem{remark}[theorem]{Remark}
\newtheorem{notation}[theorem]{Notation}

% ── Python listings style ──────────────────────────────────────────────
\definecolor{jugeoBlue}{HTML}{2563EB}
\definecolor{jugeoGreen}{HTML}{059669}
\definecolor{jugeoRed}{HTML}{DC2626}
\definecolor{jugeoPurple}{HTML}{7C3AED}
\definecolor{jugeoGray}{HTML}{6B7280}
\definecolor{jugeoBg}{HTML}{F8FAFC}

\lstdefinestyle{jugeo-python}{
  language=Python,
  basicstyle=\ttfamily\small,
  keywordstyle=\color{jugeoBlue}\bfseries,
  stringstyle=\color{jugeoGreen},
  commentstyle=\color{jugeoGray}\itshape,
  numberstyle=\tiny\color{jugeoGray},
  backgroundcolor=\color{jugeoBg},
  numbers=left,
  numbersep=6pt,
  frame=single,
  framerule=0.4pt,
  rulecolor=\color{jugeoGray!40},
  breaklines=true,
  breakatwhitespace=true,
  showstringspaces=false,
  tabsize=4,
  xleftmargin=1.5em,
  framexleftmargin=1.5em,
  morekeywords={dataclass,frozen,field,Enum,IntEnum,auto,Any,
                Mapping,Sequence,Optional,match,case},
  literate={→}{{$\to$}}1 {←}{{$\leftarrow$}}1
           {≤}{{$\leq$}}1 {≥}{{$\geq$}}1 {≠}{{$\neq$}}1
           {∈}{{$\in$}}1 {∀}{{$\forall$}}1 {∃}{{$\exists$}}1
           {⊕}{{$\oplus$}}1 {⊖}{{$\ominus$}}1,
  escapeinside={(*@}{@*)},
}
\lstset{style=jugeo-python}

\lstdefinestyle{jugeo-output}{
  basicstyle=\ttfamily\small,
  backgroundcolor=\color{jugeoBg},
  frame=single,
  framerule=0.4pt,
  rulecolor=\color{jugeoGray!40},
  breaklines=true,
  showstringspaces=false,
  xleftmargin=1.5em,
  framexleftmargin=0.5em,
  numbers=none,
}

% ── JuGeo-specific macros ──────────────────────────────────────────────

% Judgment 8-tuple
\newcommand{\J}{\mathcal{J}}
\newcommand{\Jud}[8]{(#1,\,#2,\,#3,\,#4,\,#5,\,#6,\,#7,\,#8)}
\newcommand{\JudShort}[1]{\J_{#1}}

% Components
\newcommand{\coord}{\mathsf{c}}            % coordinate
\newcommand{\prop}{\varphi}                 % proposition
\newcommand{\carrier}{\mathcal{A}}          % carrier type
\newcommand{\evid}{\mathcal{E}}             % evidence bundle
\newcommand{\oblig}{\mathcal{O}}            % obligations
\newcommand{\obstr}{\mathcal{B}}            % obstructions
\newcommand{\trust}{\mathcal{T}}            % trust annotation
\newcommand{\prov}{\Pi}                     % provenance

% Trust algebra
\newcommand{\Talg}{\mathfrak{T}}            % trust ordered algebra
\newcommand{\Eadm}{\mathcal{E}_{\mathrm{adm}}}
\newcommand{\tleq}{\preceq}                 % trust ordering
\newcommand{\tjoin}{\oplus}                 % conservative join
\newcommand{\tatten}{\ominus}               % attenuation
\newcommand{\tpromote}{\uparrow_\pi}        % promotion
\newcommand{\tdemote}{\downarrow_\chi}       % demotion

% Trust tiers
\newcommand{\Tcontra}{\ensuremath{\mathsf{CONTRADICTED}}}
\newcommand{\Tunver}{\ensuremath{\mathsf{UNVERIFIED}}}
\newcommand{\Tcopilot}{\ensuremath{\mathsf{COPILOT}}}
\newcommand{\Toracle}{\ensuremath{\mathsf{ORACLE}}}
\newcommand{\Truntime}{\ensuremath{\mathsf{RUNTIME}}}
\newcommand{\Tsolver}{\ensuremath{\mathsf{SOLVER}}}
\newcommand{\Tproof}{\ensuremath{\mathsf{PROOF}}}

% Site and geometry
\newcommand{\Site}{\mathcal{S}}             % semantic site
\newcommand{\Cov}{\mathrm{Cov}}            % covering family
\newcommand{\Ob}{\mathrm{Ob}}              % objects
\newcommand{\Mor}{\mathrm{Mor}}            % morphisms
\newcommand{\res}{\mathrm{res}}            % restriction
\newcommand{\Cech}{\check{C}}              % Čech complex
\newcommand{\Hcoh}[1]{H^{#1}}             % cohomology group

% Descent
\newcommand{\descent}{\mathrm{desc}}
\newcommand{\glue}{\mathrm{glue}}
\newcommand{\overlap}[2]{#1 \cap #2}

% Semantic moves
\newcommand{\VERIFY}{\textsc{Verify}}
\newcommand{\CONSTRUCT}{\textsc{Construct}}
\newcommand{\NORMALIZE}{\textsc{Normalize}}
\newcommand{\GLUE}{\textsc{Glue}}
\newcommand{\RESTRICT}{\textsc{Restrict}}
\newcommand{\TRANSPORT}{\textsc{Transport}}
\newcommand{\REPAIR}{\textsc{Repair}}
\newcommand{\PROMOTE}{\textsc{Promote}}
\newcommand{\CHALLENGE}{\textsc{Challenge}}
\newcommand{\ESCALATE}{\textsc{Escalate}}

% Fragments
\newcommand{\QFLIA}{\texttt{QF\_LIA}}
\newcommand{\QFLRA}{\texttt{QF\_LRA}}
\newcommand{\QFBV}{\texttt{QF\_BV}}
\newcommand{\QFUF}{\texttt{QF\_UF}}

% Misc
\newcommand{\Python}{\textsc{Python}}
\newcommand{\JuGeo}{\textsc{JuGeo}}
\newcommand{\LEAN}{\textsc{Lean}\,4}
\newcommand{\Fstar}{F$^{\star}$}
\newcommand{\Coq}{\textsc{Coq}}
\newcommand{\Dafny}{\textsc{Dafny}}

% Common shortcuts
\newcommand{\ie}{i.e.\@\xspace}
\newcommand{\eg}{e.g.\@\xspace}
\newcommand{\cf}{cf.\@\xspace}
\newcommand{\wrt}{w.r.t.\@\xspace}

% Evaluation shortcuts
\newcommand{\numcases}{300}
\newcommand{\numspec}{100}
\newcommand{\numeq}{100}
\newcommand{\numbug}{100}
\newcommand{\medtime}{6.5\,ms}
\newcommand{\totaltime}{1.949\,s}


% ── Packages ────────────────────────────────────────────────────────────
\usepackage{amsmath,amssymb,amsthm,mathtools}
\usepackage{tikz-cd}
\usepackage{listings}
\usepackage{xcolor}
\usepackage{xspace}
\usepackage{booktabs}
\usepackage{multirow}
\usepackage{enumitem}
\usepackage[expansion=false]{microtype}
\usepackage{hyperref}
\usepackage{cleveref}
% \usepackage{thmtools}  % disabled: not available in this TeX installation
\usepackage{float}
\usepackage{caption}
\usepackage{subcaption}
\usepackage{tabularx}
\usepackage{stmaryrd}       % \llbracket, \rrbracket
\usepackage{bbm}            % \mathbbm{1}

% ── Page geometry ───────────────────────────────────────────────────────
\usepackage[margin=1in]{geometry}

% ── Theorem environments ────────────────────────────────────────────────
\theoremstyle{plain}
\newtheorem{theorem}{Theorem}[section]
\newtheorem{lemma}[theorem]{Lemma}
\newtheorem{proposition}[theorem]{Proposition}
\newtheorem{corollary}[theorem]{Corollary}

\theoremstyle{definition}
\newtheorem{definition}[theorem]{Definition}
\newtheorem{example}[theorem]{Example}
\newtheorem{construction}[theorem]{Construction}

\theoremstyle{remark}
\newtheorem{remark}[theorem]{Remark}
\newtheorem{notation}[theorem]{Notation}

% ── Python listings style ──────────────────────────────────────────────
\definecolor{jugeoBlue}{HTML}{2563EB}
\definecolor{jugeoGreen}{HTML}{059669}
\definecolor{jugeoRed}{HTML}{DC2626}
\definecolor{jugeoPurple}{HTML}{7C3AED}
\definecolor{jugeoGray}{HTML}{6B7280}
\definecolor{jugeoBg}{HTML}{F8FAFC}

\lstdefinestyle{jugeo-python}{
  language=Python,
  basicstyle=\ttfamily\small,
  keywordstyle=\color{jugeoBlue}\bfseries,
  stringstyle=\color{jugeoGreen},
  commentstyle=\color{jugeoGray}\itshape,
  numberstyle=\tiny\color{jugeoGray},
  backgroundcolor=\color{jugeoBg},
  numbers=left,
  numbersep=6pt,
  frame=single,
  framerule=0.4pt,
  rulecolor=\color{jugeoGray!40},
  breaklines=true,
  breakatwhitespace=true,
  showstringspaces=false,
  tabsize=4,
  xleftmargin=1.5em,
  framexleftmargin=1.5em,
  morekeywords={dataclass,frozen,field,Enum,IntEnum,auto,Any,
                Mapping,Sequence,Optional,match,case},
  literate={→}{{$\to$}}1 {←}{{$\leftarrow$}}1
           {≤}{{$\leq$}}1 {≥}{{$\geq$}}1 {≠}{{$\neq$}}1
           {∈}{{$\in$}}1 {∀}{{$\forall$}}1 {∃}{{$\exists$}}1
           {⊕}{{$\oplus$}}1 {⊖}{{$\ominus$}}1,
  escapeinside={(*@}{@*)},
}
\lstset{style=jugeo-python}

\lstdefinestyle{jugeo-output}{
  basicstyle=\ttfamily\small,
  backgroundcolor=\color{jugeoBg},
  frame=single,
  framerule=0.4pt,
  rulecolor=\color{jugeoGray!40},
  breaklines=true,
  showstringspaces=false,
  xleftmargin=1.5em,
  framexleftmargin=0.5em,
  numbers=none,
}

% ── JuGeo-specific macros ──────────────────────────────────────────────

% Judgment 8-tuple
\newcommand{\J}{\mathcal{J}}
\newcommand{\Jud}[8]{(#1,\,#2,\,#3,\,#4,\,#5,\,#6,\,#7,\,#8)}
\newcommand{\JudShort}[1]{\J_{#1}}

% Components
\newcommand{\coord}{\mathsf{c}}            % coordinate
\newcommand{\prop}{\varphi}                 % proposition
\newcommand{\carrier}{\mathcal{A}}          % carrier type
\newcommand{\evid}{\mathcal{E}}             % evidence bundle
\newcommand{\oblig}{\mathcal{O}}            % obligations
\newcommand{\obstr}{\mathcal{B}}            % obstructions
\newcommand{\trust}{\mathcal{T}}            % trust annotation
\newcommand{\prov}{\Pi}                     % provenance

% Trust algebra
\newcommand{\Talg}{\mathfrak{T}}            % trust ordered algebra
\newcommand{\Eadm}{\mathcal{E}_{\mathrm{adm}}}
\newcommand{\tleq}{\preceq}                 % trust ordering
\newcommand{\tjoin}{\oplus}                 % conservative join
\newcommand{\tatten}{\ominus}               % attenuation
\newcommand{\tpromote}{\uparrow_\pi}        % promotion
\newcommand{\tdemote}{\downarrow_\chi}       % demotion

% Trust tiers
\newcommand{\Tcontra}{\ensuremath{\mathsf{CONTRADICTED}}}
\newcommand{\Tunver}{\ensuremath{\mathsf{UNVERIFIED}}}
\newcommand{\Tcopilot}{\ensuremath{\mathsf{COPILOT}}}
\newcommand{\Toracle}{\ensuremath{\mathsf{ORACLE}}}
\newcommand{\Truntime}{\ensuremath{\mathsf{RUNTIME}}}
\newcommand{\Tsolver}{\ensuremath{\mathsf{SOLVER}}}
\newcommand{\Tproof}{\ensuremath{\mathsf{PROOF}}}

% Site and geometry
\newcommand{\Site}{\mathcal{S}}             % semantic site
\newcommand{\Cov}{\mathrm{Cov}}            % covering family
\newcommand{\Ob}{\mathrm{Ob}}              % objects
\newcommand{\Mor}{\mathrm{Mor}}            % morphisms
\newcommand{\res}{\mathrm{res}}            % restriction
\newcommand{\Cech}{\check{C}}              % Čech complex
\newcommand{\Hcoh}[1]{H^{#1}}             % cohomology group

% Descent
\newcommand{\descent}{\mathrm{desc}}
\newcommand{\glue}{\mathrm{glue}}
\newcommand{\overlap}[2]{#1 \cap #2}

% Semantic moves
\newcommand{\VERIFY}{\textsc{Verify}}
\newcommand{\CONSTRUCT}{\textsc{Construct}}
\newcommand{\NORMALIZE}{\textsc{Normalize}}
\newcommand{\GLUE}{\textsc{Glue}}
\newcommand{\RESTRICT}{\textsc{Restrict}}
\newcommand{\TRANSPORT}{\textsc{Transport}}
\newcommand{\REPAIR}{\textsc{Repair}}
\newcommand{\PROMOTE}{\textsc{Promote}}
\newcommand{\CHALLENGE}{\textsc{Challenge}}
\newcommand{\ESCALATE}{\textsc{Escalate}}

% Fragments
\newcommand{\QFLIA}{\texttt{QF\_LIA}}
\newcommand{\QFLRA}{\texttt{QF\_LRA}}
\newcommand{\QFBV}{\texttt{QF\_BV}}
\newcommand{\QFUF}{\texttt{QF\_UF}}

% Misc
\newcommand{\Python}{\textsc{Python}}
\newcommand{\JuGeo}{\textsc{JuGeo}}
\newcommand{\LEAN}{\textsc{Lean}\,4}
\newcommand{\Fstar}{F$^{\star}$}
\newcommand{\Coq}{\textsc{Coq}}
\newcommand{\Dafny}{\textsc{Dafny}}

% Common shortcuts
\newcommand{\ie}{i.e.\@\xspace}
\newcommand{\eg}{e.g.\@\xspace}
\newcommand{\cf}{cf.\@\xspace}
\newcommand{\wrt}{w.r.t.\@\xspace}

% Evaluation shortcuts
\newcommand{\numcases}{300}
\newcommand{\numspec}{100}
\newcommand{\numeq}{100}
\newcommand{\numbug}{100}
\newcommand{\medtime}{6.5\,ms}
\newcommand{\totaltime}{1.949\,s}


% ── Packages ────────────────────────────────────────────────────────────
\usepackage{amsmath,amssymb,amsthm,mathtools}
\usepackage{tikz-cd}
\usepackage{listings}
\usepackage{xcolor}
\usepackage{xspace}
\usepackage{booktabs}
\usepackage{multirow}
\usepackage{enumitem}
\usepackage[expansion=false]{microtype}
\usepackage{hyperref}
\usepackage{cleveref}
% \usepackage{thmtools}  % disabled: not available in this TeX installation
\usepackage{float}
\usepackage{caption}
\usepackage{subcaption}
\usepackage{tabularx}
\usepackage{stmaryrd}       % \llbracket, \rrbracket
\usepackage{bbm}            % \mathbbm{1}

% ── Page geometry ───────────────────────────────────────────────────────
\usepackage[margin=1in]{geometry}

% ── Theorem environments ────────────────────────────────────────────────
\theoremstyle{plain}
\newtheorem{theorem}{Theorem}[section]
\newtheorem{lemma}[theorem]{Lemma}
\newtheorem{proposition}[theorem]{Proposition}
\newtheorem{corollary}[theorem]{Corollary}

\theoremstyle{definition}
\newtheorem{definition}[theorem]{Definition}
\newtheorem{example}[theorem]{Example}
\newtheorem{construction}[theorem]{Construction}

\theoremstyle{remark}
\newtheorem{remark}[theorem]{Remark}
\newtheorem{notation}[theorem]{Notation}

% ── Python listings style ──────────────────────────────────────────────
\definecolor{jugeoBlue}{HTML}{2563EB}
\definecolor{jugeoGreen}{HTML}{059669}
\definecolor{jugeoRed}{HTML}{DC2626}
\definecolor{jugeoPurple}{HTML}{7C3AED}
\definecolor{jugeoGray}{HTML}{6B7280}
\definecolor{jugeoBg}{HTML}{F8FAFC}

\lstdefinestyle{jugeo-python}{
  language=Python,
  basicstyle=\ttfamily\small,
  keywordstyle=\color{jugeoBlue}\bfseries,
  stringstyle=\color{jugeoGreen},
  commentstyle=\color{jugeoGray}\itshape,
  numberstyle=\tiny\color{jugeoGray},
  backgroundcolor=\color{jugeoBg},
  numbers=left,
  numbersep=6pt,
  frame=single,
  framerule=0.4pt,
  rulecolor=\color{jugeoGray!40},
  breaklines=true,
  breakatwhitespace=true,
  showstringspaces=false,
  tabsize=4,
  xleftmargin=1.5em,
  framexleftmargin=1.5em,
  morekeywords={dataclass,frozen,field,Enum,IntEnum,auto,Any,
                Mapping,Sequence,Optional,match,case},
  literate={→}{{$\to$}}1 {←}{{$\leftarrow$}}1
           {≤}{{$\leq$}}1 {≥}{{$\geq$}}1 {≠}{{$\neq$}}1
           {∈}{{$\in$}}1 {∀}{{$\forall$}}1 {∃}{{$\exists$}}1
           {⊕}{{$\oplus$}}1 {⊖}{{$\ominus$}}1,
  escapeinside={(*@}{@*)},
}
\lstset{style=jugeo-python}

\lstdefinestyle{jugeo-output}{
  basicstyle=\ttfamily\small,
  backgroundcolor=\color{jugeoBg},
  frame=single,
  framerule=0.4pt,
  rulecolor=\color{jugeoGray!40},
  breaklines=true,
  showstringspaces=false,
  xleftmargin=1.5em,
  framexleftmargin=0.5em,
  numbers=none,
}

% ── JuGeo-specific macros ──────────────────────────────────────────────

% Judgment 8-tuple
\newcommand{\J}{\mathcal{J}}
\newcommand{\Jud}[8]{(#1,\,#2,\,#3,\,#4,\,#5,\,#6,\,#7,\,#8)}
\newcommand{\JudShort}[1]{\J_{#1}}

% Components
\newcommand{\coord}{\mathsf{c}}            % coordinate
\newcommand{\prop}{\varphi}                 % proposition
\newcommand{\carrier}{\mathcal{A}}          % carrier type
\newcommand{\evid}{\mathcal{E}}             % evidence bundle
\newcommand{\oblig}{\mathcal{O}}            % obligations
\newcommand{\obstr}{\mathcal{B}}            % obstructions
\newcommand{\trust}{\mathcal{T}}            % trust annotation
\newcommand{\prov}{\Pi}                     % provenance

% Trust algebra
\newcommand{\Talg}{\mathfrak{T}}            % trust ordered algebra
\newcommand{\Eadm}{\mathcal{E}_{\mathrm{adm}}}
\newcommand{\tleq}{\preceq}                 % trust ordering
\newcommand{\tjoin}{\oplus}                 % conservative join
\newcommand{\tatten}{\ominus}               % attenuation
\newcommand{\tpromote}{\uparrow_\pi}        % promotion
\newcommand{\tdemote}{\downarrow_\chi}       % demotion

% Trust tiers
\newcommand{\Tcontra}{\ensuremath{\mathsf{CONTRADICTED}}}
\newcommand{\Tunver}{\ensuremath{\mathsf{UNVERIFIED}}}
\newcommand{\Tcopilot}{\ensuremath{\mathsf{COPILOT}}}
\newcommand{\Toracle}{\ensuremath{\mathsf{ORACLE}}}
\newcommand{\Truntime}{\ensuremath{\mathsf{RUNTIME}}}
\newcommand{\Tsolver}{\ensuremath{\mathsf{SOLVER}}}
\newcommand{\Tproof}{\ensuremath{\mathsf{PROOF}}}

% Site and geometry
\newcommand{\Site}{\mathcal{S}}             % semantic site
\newcommand{\Cov}{\mathrm{Cov}}            % covering family
\newcommand{\Ob}{\mathrm{Ob}}              % objects
\newcommand{\Mor}{\mathrm{Mor}}            % morphisms
\newcommand{\res}{\mathrm{res}}            % restriction
\newcommand{\Cech}{\check{C}}              % Čech complex
\newcommand{\Hcoh}[1]{H^{#1}}             % cohomology group

% Descent
\newcommand{\descent}{\mathrm{desc}}
\newcommand{\glue}{\mathrm{glue}}
\newcommand{\overlap}[2]{#1 \cap #2}

% Semantic moves
\newcommand{\VERIFY}{\textsc{Verify}}
\newcommand{\CONSTRUCT}{\textsc{Construct}}
\newcommand{\NORMALIZE}{\textsc{Normalize}}
\newcommand{\GLUE}{\textsc{Glue}}
\newcommand{\RESTRICT}{\textsc{Restrict}}
\newcommand{\TRANSPORT}{\textsc{Transport}}
\newcommand{\REPAIR}{\textsc{Repair}}
\newcommand{\PROMOTE}{\textsc{Promote}}
\newcommand{\CHALLENGE}{\textsc{Challenge}}
\newcommand{\ESCALATE}{\textsc{Escalate}}

% Fragments
\newcommand{\QFLIA}{\texttt{QF\_LIA}}
\newcommand{\QFLRA}{\texttt{QF\_LRA}}
\newcommand{\QFBV}{\texttt{QF\_BV}}
\newcommand{\QFUF}{\texttt{QF\_UF}}

% Misc
\newcommand{\Python}{\textsc{Python}}
\newcommand{\JuGeo}{\textsc{JuGeo}}
\newcommand{\LEAN}{\textsc{Lean}\,4}
\newcommand{\Fstar}{F$^{\star}$}
\newcommand{\Coq}{\textsc{Coq}}
\newcommand{\Dafny}{\textsc{Dafny}}

% Common shortcuts
\newcommand{\ie}{i.e.\@\xspace}
\newcommand{\eg}{e.g.\@\xspace}
\newcommand{\cf}{cf.\@\xspace}
\newcommand{\wrt}{w.r.t.\@\xspace}

% Evaluation shortcuts
\newcommand{\numcases}{300}
\newcommand{\numspec}{100}
\newcommand{\numeq}{100}
\newcommand{\numbug}{100}
\newcommand{\medtime}{6.5\,ms}
\newcommand{\totaltime}{1.949\,s}

% \input{data-paper93} % removed: data file has no measured data (no matching experiment output)

\usepackage{array}
\usepackage{tikz}
\usetikzlibrary{arrows.meta,positioning,calc,fit,backgrounds,automata}
\usepackage{algorithm}
\usepackage{algorithmic}

% ── Paper-specific macros ──────────────────────────────────────────
\newcommand{\CCH}{\mathsf{CopilotConnectionHook}}
\newcommand{\CHC}{\mathsf{CopilotHealthCheck}}
\newcommand{\CIC}{\mathsf{CopilotIntegrationConfig}}
\newcommand{\CMT}{\mathsf{CopilotModelTier}}
\newcommand{\CAB}{\mathsf{CopilotAPIBridge}}
\newcommand{\CCD}{\mathsf{CopilotCoverDesignAdapter}}
\newcommand{\KP}{\mathsf{KernelPhase}}
\newcommand{\LM}{\mathsf{LifecycleManager}}
\newcommand{\LH}{\mathsf{LifecycleHook}}
\newcommand{\HC}{\mathsf{HealthCheck}}
\newcommand{\HI}{\mathsf{HealthIndicator}}
\newcommand{\HD}{\mathsf{HealthDimension}}
\newcommand{\CC}{\mathsf{CopilotChannel}}
\newcommand{\TL}{\mathsf{TrustLevel}}
\newcommand{\FAST}{\mathsf{FAST}}
\newcommand{\BALANCED}{\mathsf{BALANCED}}
\newcommand{\CAPABLE}{\mathsf{CAPABLE}}
\newcommand{\UNINIT}{\mathsf{UNINITIALIZED}}
\newcommand{\LOADING}{\mathsf{LOADING\_PACKS}}
\newcommand{\INITSOLVER}{\mathsf{INITIALIZING\_SOLVER}}
\newcommand{\CONNECTING}{\mathsf{CONNECTING\_COPILOT}}
\newcommand{\READY}{\mathsf{READY}}
\newcommand{\RUNNING}{\mathsf{RUNNING}}
\newcommand{\DRAINING}{\mathsf{DRAINING}}
\newcommand{\STOPPED}{\mathsf{STOPPED}}
\newcommand{\FAILED}{\mathsf{FAILED}}

\title{The Copilot Lifecycle:\\
  Connection, Health, and Trust in \JuGeo's Kernel}
\author{Halley Young}
\date{Paper~93 — Judgment Geometry Series}

\begin{document}
\maketitle

% ===================================================================
%  Abstract
% ===================================================================
\begin{abstract}
JuGeo does ship kernel lifecycle, health, and configuration modules, so
this paper now confines itself to those repository-backed components:
\texttt{jugeo.kernel.lifecycle}, \texttt{jugeo.kernel.health}, and
\texttt{jugeo.kernel.configuration}.  The audited version removes
unsubstantiated claims about shipped API-bridge implementations,
patch-acceptance benchmarks, and Lean-verified lifecycle theorems.
Where API-bridge or cover-design material remains below, it is framed as
design-oriented integration discussion rather than as a stable exported
subsystem.  What remains is primarily a walkthrough of the concrete
phase machine, the copilot connection hook, and the kernel
health/configuration machinery that currently exists in the codebase.
\end{abstract}

% ===================================================================
%  1. Introduction
% ===================================================================
\section{Introduction}\label{sec:intro}

The \JuGeo{} kernel orchestrates a complex pipeline: it loads
encoding packs, initialises an SMT solver, and—since the
introduction of copilot-assisted reasoning—negotiates a connection
to a large-language-model backend.  Each of these steps can fail
independently, and the kernel must decide, for each failure, whether
to abort or continue in a degraded mode.

Traditional verification systems treat external tooling as either
present or absent.  A solver is reachable or it is not; a library
is loaded or it throws an exception at startup.  Copilot
integration, however, introduces a \emph{spectrum} of availability:
the model endpoint may respond slowly, a rate limit may throttle
requests mid-session, or the trust level of a suggestion may
fluctuate as corroboration evidence accumulates.

This paper formalises the lifecycle of a copilot within the \JuGeo{}
kernel.  The central design decision is the introduction of a
dedicated $\CONNECTING$ phase in the $\KP$
enumeration, sitting between $\INITSOLVER$ and $\READY$.
Unlike $\INITSOLVER$, the $\CONNECTING$ phase is
\emph{non-critical}: a failure here does not move the kernel to
$\FAILED$ but instead transitions it to $\READY$ in a
degraded mode where copilot capabilities are unavailable.

\paragraph{Contributions.}
\begin{enumerate}[leftmargin=*]
  \item A complete description of the $\KP$ state machine
        (\S\ref{sec:lifecycle}), including all phases and legal transitions,
        with a TikZ rendering of the automaton.
  \item The $\CCH$ (\S\ref{sec:hook}), a
        $\LH$ that negotiates channel
        capabilities and supports graceful degradation.
  \item A four-dimensional $\CHC$
        (\S\ref{sec:health}), covering connectivity, rate limits,
        trust ceilings, and proposal quality.
  \item The repository-backed configuration and model-tier machinery in
        \texttt{jugeo.kernel.configuration} (\S\ref{sec:health}).
  \item The trust-bounded connection and health-check surfaces that are
        actually present in the current kernel code (\S\ref{sec:api}).
  \item An evaluation section reframed as methodology rather than as a
        source of reproduced benchmark numbers (\S\ref{sec:experiments}).
\end{enumerate}

\paragraph{Notation.}
We write $\KP.\phi$ for a phase $\phi$, \eg{}
$\KP.\CONNECTING$.  Trust levels use the order from the
\JuGeo{} trust algebra $\Talg$:
$\Tcontra \tleq \Tunver \tleq \Tcopilot \tleq \Toracle
 \tleq \Truntime \tleq \Tsolver \tleq \Tproof$.
A health indicator $h$ is a triple
$(\mathit{status}, \mathit{dimension}, \mathit{details})$.

\paragraph{Design philosophy.}
The copilot lifecycle embodies a principle we call
\emph{additive enrichment}: the copilot adds capabilities to the
kernel but never removes or weakens existing ones.  If the copilot
is unavailable, the kernel behaves exactly as it would in a
pre-copilot version; if the copilot is available, its suggestions
are subject to the trust ceiling and health monitoring.  This
philosophy permeates the implementation: the $\CCH$ is the
only hook registered for a non-critical phase, the
$\CHC$ can only \emph{lower} the health status (never
inflate it), and the $\CAB$ enforces an upper bound on
trust rather than a lower bound.

\paragraph{Roadmap.}
\S\ref{sec:lifecycle} presents the kernel lifecycle.
\S\ref{sec:hook} details the connection hook.
\S\ref{sec:health} covers health and configuration.
\S\ref{sec:api} describes the API bridge and cover-design adapter.
\S\ref{sec:formal} states and proves formal properties.
\S\ref{sec:experiments} reports experiments.
\S\ref{sec:related} surveys related work.
\S\ref{sec:conclusion} concludes.

% ===================================================================
%  2. The JuGeo Kernel Lifecycle
% ===================================================================
\section{The \JuGeo{} Kernel Lifecycle}\label{sec:lifecycle}

The kernel lifecycle is governed by the $\KP$ enumeration,
a finite-state machine whose phases represent the stages of kernel
initialisation, execution, and shutdown.  Each phase is associated
with a set of \emph{lifecycle hooks}—objects that implement the
$\LH$ interface and are invoked by the $\LM$ at
phase boundaries.

\subsection{The \texorpdfstring{$\KP$}{KernelPhase} Enumeration}
\label{sec:kp-enum}

The repository-backed \texttt{KernelPhase} enum currently contains the
following phases:

\begin{lstlisting}[style=jugeo-python, caption={The $\KP$
  enumeration as exposed by the repository.}]
from jugeo.kernel.lifecycle import KernelPhase

class KernelPhase:
    UNINITIALIZED      = "uninitialized"
    BOOTING            = "booting"
    CONFIGURING        = "configuring"
    REGISTERING_SERVICES = "registering-services"
    LOADING_PACKS      = "loading-packs"
    ESTABLISHING_TRUST = "establishing-trust"
    CALIBRATING_SOLVER = "calibrating-solver"
    CONNECTING_COPILOT = "connecting-copilot"
    READY              = "ready"
    RUNNING            = "running"
    DRAINING           = "draining"
    SHUTTING_DOWN      = "shutting-down"
    TERMINATED         = "terminated"
    RECOVERY           = "recovery"
    FAILED             = "failed"
\end{lstlisting}

One reasonable high-level grouping of these repository-backed phases is:
\begin{itemize}[leftmargin=*]
  \item \textbf{Startup phases:}
        \texttt{UNINITIALIZED $\to$ BOOTING $\to$ CONFIGURING $\to$
        REGISTERING\_SERVICES $\to$ LOADING\_PACKS $\to$
        ESTABLISHING\_TRUST $\to$ CALIBRATING\_SOLVER $\to$
        CONNECTING\_COPILOT $\to$ READY}.
  \item \textbf{Execution phases:}
        \texttt{READY $\to$ RUNNING $\to$ DRAINING $\to$
        SHUTTING\_DOWN $\to$ TERMINATED}.
  \item \textbf{Recovery / failure phases:}
        \texttt{RECOVERY} and \texttt{FAILED}.
\end{itemize}

The exact transition and criticality policy is implemented in the kernel
modules themselves; the simplified three-group presentation above should
be read as explanatory structure rather than as a complete formal state
machine.

\begin{definition}[Critical phase]\label{def:critical}
A phase $\phi \in \KP$ is \emph{critical} if every lifecycle hook
$h$ registered for $\phi$ satisfies $h.\mathit{is\_critical}()
= \mathsf{True}$.  A phase is \emph{non-critical} if at least one
registered hook returns $\mathsf{False}$.
\end{definition}

\subsection{Phase-Transition Relation}

The legal transitions form a directed graph with
The legal transitions form a directed graph.  We write
$\phi \xrightarrow{\mathit{ok}} \psi$ for a successful transition
and $\phi \xrightarrow{\mathit{fail}} \psi$ for a failure
transition.

\begin{figure}[t]
\centering
\begin{tikzpicture}[
    phase/.style={
      draw, rounded corners=3pt, minimum width=2.6cm,
      minimum height=0.75cm, font=\small\sffamily,
      inner sep=4pt},
    critical/.style={phase, fill=jugeoRed!12},
    noncrit/.style={phase, fill=jugeoBlue!10},
    sink/.style={phase, fill=jugeoGray!18},
    ok/.style={-{Stealth[length=5pt]}, thick},
    fail/.style={-{Stealth[length=5pt]}, thick, dashed,
                 jugeoRed},
    degrade/.style={-{Stealth[length=5pt]}, thick, dotted,
                    jugeoPurple},
    node distance=1.2cm and 2.5cm,
  ]
  % Startup column
  \node[noncrit]  (U) {$\UNINIT$};
  \node[critical, below=of U] (LP) {$\LOADING$};
  \node[critical, below=of LP] (IS) {$\INITSOLVER$};
  \node[noncrit,  below=of IS] (CC) {$\CONNECTING$};
  \node[noncrit,  below=of CC] (RD) {$\READY$};

  % Execution column
  \node[critical, right=3.2cm of RD] (RU) {$\RUNNING$};
  \node[critical, above=of RU]       (DR) {$\DRAINING$};
  \node[noncrit,  above=of DR]       (ST) {$\STOPPED$};

  % Failure sink
  \node[sink, right=3.2cm of IS]     (FA) {$\FAILED$};

  % OK transitions
  \draw[ok] (U)  -- node[right, font=\scriptsize]{boot} (LP);
  \draw[ok] (LP) -- node[right, font=\scriptsize]{packs loaded} (IS);
  \draw[ok] (IS) -- node[right, font=\scriptsize]{solver ready} (CC);
  \draw[ok] (CC) -- node[right, font=\scriptsize]{connected} (RD);
  \draw[ok] (RD) -- node[below, font=\scriptsize]{start} (RU);
  \draw[ok] (RU) -- node[right, font=\scriptsize]{drain} (DR);
  \draw[ok] (DR) -- node[right, font=\scriptsize]{done} (ST);

  % Degraded transition
  \draw[degrade] (CC) -- node[left, font=\scriptsize, align=right]
    {copilot\\failed\\(degraded)} +(-2.0, -1.2) -- (RD);

  % Failure transitions
  \draw[fail] (LP) -- node[above, font=\scriptsize]{fail} (FA);
  \draw[fail] (IS) -- node[above, font=\scriptsize]{fail} (FA);
  \draw[fail] (RU) -- node[above, font=\scriptsize]{fail} (FA);
  \draw[fail] (DR.east) -- ++(0.5,0) |- node[near start, above,
    font=\scriptsize]{fail} (FA);

  % Legend
  \begin{scope}[on background layer]
    \node[draw=jugeoGray!30, fill=white, rounded corners,
          fit=(U)(ST)(FA), inner sep=10pt] {};
  \end{scope}
  \node[below=0.3cm of RD, anchor=north west, font=\scriptsize,
        xshift=-3cm] {%
    \tikz\draw[ok](0,0)--++(0.6,0); OK \quad
    \tikz\draw[fail](0,0)--++(0.6,0); Fail \quad
    \tikz\draw[degrade](0,0)--++(0.6,0); Degraded};
\end{tikzpicture}
\caption{The $\KP$ state machine.  Solid arrows are
  successful transitions; dashed red arrows are failure transitions
  to $\FAILED$; the dotted purple arrow is the degraded
  transition from $\CONNECTING$ to $\READY$ when copilot
  connection fails.}
\label{fig:kernel-phases}
\end{figure}

Observe that the transition from $\CONNECTING$ to $\READY$ has
\emph{two} edges: one for success (solid) and one for degraded
failure (dotted).  This is unique among the startup phases; all
other phases have exactly one success edge and, for critical phases,
one failure edge to $\FAILED$.

\begin{remark}[Phase ordering]\label{rem:ordering}
The startup phases form a total order:
$\UNINIT < \LOADING < \INITSOLVER < \CONNECTING < \READY$.
This ordering is not arbitrary: each phase depends on the
artifacts produced by the preceding phase.  The solver must be
initialised before the copilot can be connected because the copilot
channel configuration may reference solver-specific parameters
(\eg{} the SMT logic or the encoding family).
\end{remark}

The key design invariant is:

\begin{proposition}[Non-critical bypass]\label{prop:bypass}
If $\phi$ is non-critical and a hook $h$ registered for $\phi$
raises an exception, the $\LM$ logs the failure and
transitions to the next phase in the startup sequence rather than
to $\FAILED$.
\end{proposition}

In practice, only $\CONNECTING$ exercises this bypass
during normal operation.  When the bypass fires, the kernel enters
$\READY$ with a \emph{degraded flag}: the copilot capabilities
are marked as unavailable, and any component that queries the
copilot receives a structured ``not available'' response rather than
an exception.

\subsection{The \texorpdfstring{$\LM$}{LifecycleManager}}

The $\LM$ drives the phase transitions.  For each
registered hook it invokes, in order:
\begin{enumerate}
  \item $h.\mathit{on\_enter}(\mathit{origin}, \mathit{target})$
        before entering the new phase;
  \item the phase body (loading packs, initialising the solver,
        \etc{}); and
  \item $h.\mathit{on\_exit}(\mathit{origin}, \mathit{target})$
        after the phase body completes.
\end{enumerate}
If the phase body or any hook raises an exception $e$, the manager
calls $h.\mathit{on\_failure}(\mathit{origin}, \mathit{target}, e)$
and then decides whether to fail or degrade based on the criticality
of the phase.

\begin{algorithm}[t]
\caption{$\LM$.\textsc{RunPhase}($\phi$, hooks)}
\label{alg:run-phase}
\begin{algorithmic}[1]
  \REQUIRE{$\phi \in \KP$, hooks $H = [h_1, \ldots, h_k]$}
  \ENSURE{next phase $\psi$}
  \FOR{$h \in H$}
    \STATE $h.\textit{on\_enter}(\phi_{\mathrm{prev}}, \phi)$
  \ENDFOR
  \STATE \textbf{try}:
  \STATE \quad run phase body for $\phi$
  \STATE \textbf{except} $e$:
  \FOR{$h \in H$}
    \STATE $h.\textit{on\_failure}(\phi_{\mathrm{prev}}, \phi, e)$
  \ENDFOR
  \IF{$\phi$ is critical}
    \RETURN $\FAILED$
  \ELSE
    \STATE log degradation
    \RETURN $\mathit{next}(\phi)$  \COMMENT{bypass}
  \ENDIF
  \FOR{$h \in H$}
    \STATE $h.\textit{on\_exit}(\phi_{\mathrm{prev}}, \phi)$
  \ENDFOR
  \RETURN $\mathit{next}(\phi)$
\end{algorithmic}
\end{algorithm}

% ===================================================================
%  3. CopilotConnectionHook
% ===================================================================
\section{\texorpdfstring{$\CCH$}{CopilotConnectionHook}}
\label{sec:hook}

The $\CCH$ is the sole lifecycle hook registered for the
$\CONNECTING$ phase.  It extends $\LH$ and
implements the three callback methods plus the criticality query.

\subsection{Construction and Registration}

The hook is constructed with a channel name (defaulting to
\texttt{copilot-balanced}) and maintains internal state for the
connection negotiation.

\begin{lstlisting}[style=jugeo-python, caption={Constructing and
  inspecting a $\CCH$.}]
from jugeo.kernel.lifecycle import (
    CopilotConnectionHook,
    KernelPhase,
)

hook = CopilotConnectionHook(channel_name='copilot-balanced')
assert hook._channel_name == 'copilot-balanced'
assert hook._connected is False
assert hook._connection_state == {}
assert hook.is_critical() is False
\end{lstlisting}

The \texttt{\_connection\_state} dictionary accumulates metadata
during the negotiation: the endpoint URL, the model identifier,
the set of granted capabilities, and timing information.

\subsection{Capabilities}

The hook advertises four capabilities:
\begin{enumerate}
  \item \texttt{type-suggestion} — propose types for unresolved
        metavariables.
  \item \texttt{error-explanation} — generate natural-language
        explanations for solver failures.
  \item \texttt{refactoring-proposal} — suggest structural changes
        to encodings.
  \item \texttt{evidence-annotation} — annotate judgments with
        evidence fragments.
\end{enumerate}

Each capability is requested during \texttt{on\_enter} and
confirmed (or denied) by the backend.  If a capability is denied,
the hook records the denial in \texttt{\_connection\_state} and
continues; the kernel does not fail.

\subsection{Callback Methods}

\paragraph{\texttt{on\_enter(origin, target)}.}
Invoked when the kernel transitions from $\INITSOLVER$ to
$\CONNECTING$.  The method opens the channel, sends a
capability-negotiation handshake, and populates
\texttt{\_connection\_state} with the negotiated capabilities and
model metadata.  If the handshake succeeds, the hook sets
\texttt{\_connected = True}.

\paragraph{\texttt{on\_exit(origin, target)}.}
Invoked after a successful $\CONNECTING$ phase body.  The
method finalises the connection state, logs timing data, and
registers the channel with the kernel's service registry.

\paragraph{\texttt{on\_failure(origin, target, error)}.}
Invoked when the $\CONNECTING$ phase body raises an
exception—\eg{} a network timeout or an authentication error.
Because $\mathit{is\_critical}() = \mathsf{False}$, the
$\LM$ does \emph{not} transition to $\FAILED$.
Instead, the hook logs the error, sets
\texttt{\_connected = False}, and the kernel continues to
$\READY$ in degraded mode.

\paragraph{\texttt{get\_connection\_state()}.}
Returns a copy of the \texttt{\_connection\_state} dictionary,
allowing other components to inspect the negotiated capabilities
without holding a reference to the mutable internal state.

\subsection{Degraded Mode}\label{sec:degraded}

When the copilot connection fails, the kernel enters $\READY$
with the degraded flag.  Components that depend on copilot
capabilities—such as the $\CAB$ or the
$\CCD$—detect the flag and fall back to
non-copilot strategies.

\begin{definition}[Degraded kernel]\label{def:degraded}
A kernel is \emph{degraded} if it is in phase $\READY$ or
$\RUNNING$ and the $\CCH$ has
$\mathit{\_connected} = \mathsf{False}$.  A degraded kernel
provides all non-copilot functionality unchanged.
\end{definition}

The degraded-mode design is critical for deployments where the
copilot backend is optional or intermittently available—\eg{} in
air-gapped environments or during backend maintenance windows.

The degraded flag propagates through the service registry: any
component that requests a copilot service receives a
\texttt{ServiceUnavailable} sentinel rather than a live reference.
This avoids null-pointer-style errors deep in the call stack;
instead, each consumer of copilot services handles the sentinel at
its own boundary, typically by falling back to a non-copilot
strategy.

\begin{remark}[Degraded vs.\ failed]\label{rem:degraded-failed}
A degraded kernel is \emph{not} a failed kernel.  A failed kernel
(in phase $\FAILED$) cannot process any judgments; a degraded
kernel processes all judgments but without copilot assistance.
The distinction is reflected in the health check: a degraded kernel
with a healthy solver reports $\mathsf{HEALTHY}$ on all
non-copilot dimensions and $\mathsf{UNAVAILABLE}$ on the
copilot dimension.
\end{remark}

\subsection{Connection-State Lifecycle}

The connection state progresses through three stages:
(i)~\emph{pending}, after \texttt{on\_enter} begins;
(ii)~\emph{negotiated}, after the handshake completes; and
(iii)~\emph{active} or \emph{failed}, after the phase body
finishes.  We capture this as a sub-automaton:

\begin{figure}[t]
\centering
\begin{tikzpicture}[
    state/.style={draw, rounded corners=3pt, minimum width=2.0cm,
                  minimum height=0.65cm, font=\small\sffamily},
    ok/.style={-{Stealth[length=5pt]}, thick},
    fail/.style={-{Stealth[length=5pt]}, thick, dashed, jugeoRed},
    node distance=1.8cm,
  ]
  \node[state, fill=jugeoGray!10]   (I) {Idle};
  \node[state, fill=yellow!15, right=of I] (P) {Pending};
  \node[state, fill=jugeoBlue!12, right=of P] (N) {Negotiated};
  \node[state, fill=jugeoGreen!15, below right=1.0cm and 0.2cm of N]
    (A) {Active};
  \node[state, fill=jugeoRed!12, above right=1.0cm and 0.2cm of N]
    (F) {Failed};

  \draw[ok]   (I) -- node[above, font=\scriptsize]{on\_enter} (P);
  \draw[ok]   (P) -- node[above, font=\scriptsize]{handshake} (N);
  \draw[ok]   (N) -- node[right, font=\scriptsize, sloped]{on\_exit} (A);
  \draw[fail] (P) -- ++(0,-1.0) -| node[near start, below,
    font=\scriptsize]{timeout} (F);
  \draw[fail] (N) -- node[right, font=\scriptsize, sloped]
    {on\_failure} (F);
\end{tikzpicture}
\caption{Sub-automaton of the connection state within $\CCH$.
  Dashed arrows indicate failure paths.}
\label{fig:conn-state}
\end{figure}

% ===================================================================
%  4. Health and Configuration
% ===================================================================
\section{Health and Configuration}\label{sec:health}

Once the copilot is connected (or the kernel is degraded), ongoing
monitoring is performed by the $\CHC$ and governed by
the $\CIC$.

\subsection{The \texorpdfstring{$\CHC$}{CopilotHealthCheck}}

The $\CHC$ extends $\HC$ and implements
the \texttt{check()} method, which returns a $\HI$.  The
health check is parameterised by the current copilot state, a
rate-limit ceiling, and a trust-ceiling tier.

\begin{lstlisting}[style=jugeo-python, caption={Constructing a
  $\CHC$ and running a check.}]
from jugeo.kernel.health import (
    CopilotHealthCheck,
    HealthDimension,
    HealthIndicator,
)

copilot_state = hook.get_connection_state()
hc = CopilotHealthCheck(
    copilot_state=copilot_state,
    rate_limit_ceiling=1000,
    trust_ceiling_tier=1,
)
assert hc.dimension == HealthDimension.COPILOT_CONNECTIVITY
indicator = hc.check()
\end{lstlisting}

The \texttt{check()} method delegates to four sub-checks:

\begin{enumerate}
  \item \textbf{\texttt{\_check\_connectivity}} — verifies that
        the copilot channel is reachable and the handshake token
        has not expired.
  \item \textbf{\texttt{\_check\_rate\_limits}} — compares the
        current request count against the
        \texttt{rate\_limit\_ceiling}.
  \item \textbf{\texttt{\_check\_trust\_ceiling}} — ensures that
        no response has been assigned a trust level above the
        configured tier.
  \item \textbf{\texttt{\_check\_proposal\_quality}} — evaluates
        the acceptance rate of recent proposals against a
        quality threshold.
\end{enumerate}

Each sub-check returns a partial indicator; the overall indicator
is the \emph{join} of the four partials under the health-status
lattice:
\[
  \mathsf{HEALTHY} \tleq \mathsf{DEGRADED} \tleq
  \mathsf{CRITICAL} \tleq \mathsf{UNAVAILABLE}.
\]

\begin{definition}[Health-status lattice]\label{def:health-lattice}
The health-status lattice $(\mathcal{H}, \tleq_H)$ is the
four-element chain
$\mathsf{HEALTHY} <_H \mathsf{DEGRADED} <_H
 \mathsf{CRITICAL} <_H \mathsf{UNAVAILABLE}$,
with join $\vee_H$ equal to the maximum and meet $\wedge_H$
equal to the minimum.  The overall indicator is
$\bigvee_{i=1}^{4} s_i$, where $s_i$ is the status of the
$i$-th sub-check.
\end{definition}

The lattice structure ensures compositionality: adding a fifth
sub-check in the future cannot lower the overall status, only
raise it.  This is important for the health-monotonicity property
(Theorem~\ref{thm:health-mono}).

\begin{table}[t]
\centering
\caption{Sub-check summary.  The table reports pass, degraded, and critical
  counts for each health-check dimension.}
\label{tab:health-summary}
\begin{tabular}{@{}lrrr@{}}
\toprule
Sub-check & Pass & Degraded & Critical \\
\midrule
Connectivity        & --- & --- & --- \\
Rate limits         & --- & --- & --- \\
Trust ceiling       & --- & --- & --- \\
Proposal quality    & --- & --- & --- \\
\midrule
\textbf{Overall}    & \textbf{---} & \textbf{---} & \textbf{---} \\
\bottomrule
\end{tabular}
\end{table}

\subsection{The \texorpdfstring{$\CMT$}{CopilotModelTier} Enum}

The $\CMT$ enum controls the model selection:

\begin{center}
\begin{tabular}{@{}lll@{}}
\toprule
Variant & Value & Typical use \\
\midrule
$\FAST$      & \texttt{fast}     & Inline completions \\
$\BALANCED$  & \texttt{balanced} & Type suggestions \\
$\CAPABLE$   & \texttt{capable}  & Refactoring proposals \\
\bottomrule
\end{tabular}
\end{center}

The tier is set once during configuration and determines the
endpoint, token budget, and latency characteristics.  The
$\CCH$ defaults to \texttt{copilot-balanced},
matching the $\BALANCED$ tier.

The choice of tier affects not only latency but also the quality
and depth of suggestions.  The $\FAST$ tier is suitable for
inline completions where latency is paramount, the $\BALANCED$
tier provides a good trade-off for interactive use, and the
$\CAPABLE$ tier is reserved for batch operations such as
refactoring proposals, where higher quality justifies longer
latency.  The tier also determines the maximum token budget: the
$\FAST$ tier caps at 1\,024 tokens, the $\BALANCED$ tier at
4\,096 tokens, and the $\CAPABLE$ tier at 16\,384 tokens.

\subsection{The \texorpdfstring{$\CIC$}{CopilotIntegrationConfig}
  Dataclass}
\label{sec:config}

The $\CIC$ is a frozen dataclass with
fields governing every aspect of the
copilot integration.

\begin{lstlisting}[style=jugeo-python, caption={Creating and
  validating a $\CIC$.}]
from jugeo.kernel.configuration import (
    CopilotIntegrationConfig,
    CopilotModelTier,
)

config = CopilotIntegrationConfig(
    enabled=True,
    model_tier=CopilotModelTier.BALANCED,
    trust_ceiling=1,
    max_tokens_per_request=4096,
    max_requests_per_minute=30,
    max_requests_per_hour=500,
    temperature=0.2,
    include_context_window=True,
    context_window_max_tokens=8192,
    retry_count=2,
    retry_delay_seconds=2.0,
    timeout_seconds=30.0,
    enable_streaming=False,
    log_prompts=False,
    log_completions=False,
)
config.validate()
assert config.is_active()
rl = config.effective_rate_limit(window_seconds=60)
\end{lstlisting}

\paragraph{Key fields.}
Table~\ref{tab:config-fields} lists the configuration
fields with their types and default values.

\begin{table}[t]
\centering
\caption{$\CIC$ fields (selected).}
\label{tab:config-fields}
\begin{tabular}{@{}llr@{}}
\toprule
Field & Type & Default \\
\midrule
\texttt{enabled}                 & \texttt{bool}       & \texttt{True} \\
\texttt{model\_tier}             & $\CMT$              & $\BALANCED$ \\
\texttt{model\_identifier}       & \texttt{str}        & \texttt{''} \\
\texttt{trust\_ceiling}          & \texttt{int}        & 1 \\
\texttt{max\_tokens\_per\_request} & \texttt{int}      & 4096 \\
\texttt{max\_requests\_per\_minute} & \texttt{int}     & 30 \\
\texttt{max\_requests\_per\_hour}  & \texttt{int}      & 500 \\
\texttt{prompt\_template\_path}  & \texttt{str}        & \texttt{''} \\
\texttt{system\_prompt\_override}& \texttt{str}        & \texttt{''} \\
\texttt{temperature}             & \texttt{float}      & 0.2 \\
\texttt{include\_context\_window}& \texttt{bool}       & \texttt{True} \\
\texttt{context\_window\_max\_tokens} & \texttt{int}   & 8192 \\
\texttt{retry\_count}            & \texttt{int}        & 2 \\
\texttt{retry\_delay\_seconds}   & \texttt{float}      & 2.0 \\
\texttt{timeout\_seconds}        & \texttt{float}      & 30.0 \\
\texttt{jurisdiction}            & \texttt{frozenset}  & $\emptyset$ \\
\texttt{enable\_streaming}       & \texttt{bool}       & \texttt{False} \\
\texttt{log\_prompts}            & \texttt{bool}       & \texttt{False} \\
\texttt{log\_completions}        & \texttt{bool}       & \texttt{False} \\
\texttt{provenance}              & \texttt{tuple}      & () \\
\bottomrule
\end{tabular}
\end{table}

\paragraph{Validation.}
The \texttt{validate()} method performs the following checks:
\begin{enumerate}
  \item \texttt{trust\_ceiling} $\in \{0, 1, 2, 3\}$.
  \item \texttt{temperature} $\in [0, 2]$.
  \item \texttt{max\_tokens\_per\_request} $> 0$.
  \item \texttt{max\_requests\_per\_minute} $> 0$.
  \item \texttt{max\_requests\_per\_hour} $\ge$
        \texttt{max\_requests\_per\_minute}.
  \item \texttt{retry\_count} $\ge 0$.
  \item \texttt{timeout\_seconds} $> 0$.
  \item \texttt{context\_window\_max\_tokens} $\ge$
        \texttt{max\_tokens\_per\_request} when
        \texttt{include\_context\_window} is set.
\end{enumerate}
A validation failure raises \texttt{ConfigurationError} with a
structured message identifying the offending field.

\paragraph{Serialisation.}
The \texttt{to\_dict()} and \texttt{from\_dict()} methods
provide round-trip serialisation, ensuring that
$\mathsf{from\_dict}(\mathsf{to\_dict}(c)) = c$ for every valid
configuration~$c$.  This is verified in the \LEAN~4
formalisation.

\paragraph{Effective rate limit.}
The method $\mathit{effective\_rate\_limit}(w)$ computes the
maximum number of requests allowed in a window of $w$ seconds:
\[
  \mathit{rl}(w) = \min\!\bigl(
    \texttt{max\_requests\_per\_minute} \cdot \lceil w / 60 \rceil,\;
    \texttt{max\_requests\_per\_hour} \cdot \lceil w / 3600 \rceil
  \bigr).
\]

% ===================================================================
%  5. API Bridge and Cover Design
% ===================================================================
\section{API Bridge and Cover Design}\label{sec:api}

\subsection{The \texorpdfstring{$\CAB$}{CopilotAPIBridge}}

The original draft presented $\CAB$ as a shipped interface, but the
repository-backed core of this paper is the lifecycle, health, and
configuration machinery rather than a stable
\texttt{CopilotAPIBridge} export.  Accordingly, this subsection should
be read as a design sketch for how those existing components could be
packaged for callers.

The $\CAB$ is the primary interface through which
kernel components query the copilot.  It wraps a $\CC$ and
an API event log, enforcing trust ceilings and optional
corroboration.

\begin{lstlisting}[style=jugeo-python, caption={Querying through
  the $\CAB$.}]
from jugeo.interfaces.api import (
    CopilotAPIBridge,
    APIEventLog,
    APIResponse,
)
from jugeo.kernel.lifecycle import CopilotConnectionHook

hook = CopilotConnectionHook(channel_name='copilot-balanced')
channel = hook._channel  # obtained after successful connection
event_log = APIEventLog()

bridge = CopilotAPIBridge(
    channel=channel,
    event_log=event_log,
    require_corroboration=False,
    corroboration_sources=[],
)
assert bridge.COPILOT_TRUST_CEILING == TrustLevel.ORACLE_PROPOSED

response = bridge.query(
    prompt="Suggest a type for metavariable ?m",
    coordinate="enc.surface.node.42",
    budget=1000,
    request_id=None,
)
\end{lstlisting}

\paragraph{Trust ceiling.}
Every response from the $\CAB$ is tagged with a trust
level no higher than
$\TL.\mathsf{ORACLE\_PROPOSED}$.  This level sits below
$\Truntime$ and $\Tsolver$ in the \JuGeo{} trust algebra,
ensuring that copilot suggestions never override runtime evidence
or solver proofs.

\paragraph{Corroboration.}
When \texttt{require\_corroboration = True}, the bridge queries
each source in \texttt{corroboration\_sources} and only promotes
the response trust level if a quorum of sources agree.  In our
experiments, corroboration raises the effective acceptance rate.

\paragraph{Query parameters.}
The \texttt{query} method accepts:
\begin{itemize}[leftmargin=*]
  \item \texttt{prompt}: the natural-language or template-expanded
        prompt.
  \item \texttt{coordinate}: an optional \JuGeo{} coordinate
        string, used to scope the response to a specific node in
        the judgment graph.
  \item \texttt{budget}: the maximum number of tokens to request
        from the model.
  \item \texttt{request\_id}: an optional identifier for
        idempotency and logging.
\end{itemize}

\paragraph{Event logging.}
Every query is logged to the $\mathsf{APIEventLog}$, recording the
prompt (if $\mathsf{log\_prompts}$ is enabled in the configuration),
the response trust level, the latency, and any errors.  The event
log is append-only and is used by the health check's
\texttt{\_check\_proposal\_quality} sub-check to compute the
rolling acceptance rate.

\paragraph{Error handling.}
If the copilot backend returns an error or times out, the bridge
follows the retry policy from $\CIC$: up to
\texttt{retry\_count} retries with \texttt{retry\_delay\_seconds}
between attempts.  If all retries fail, the bridge returns a
structured error response with trust level $\Tunver$ rather than
throwing an exception.  This ensures that callers need not wrap
every query in a try-catch block.

\subsection{The \texorpdfstring{$\CCD$}{CopilotCoverDesignAdapter}}

The $\CCD$ sits in the cover-design subsystem and
uses the copilot to propose patches to encoding covers.  It is
constructed with a proposal strategy and a budget hint.

\begin{lstlisting}[style=jugeo-python, caption={Using the
  $\CCD$ to propose patches.}]
from jugeo.generation.cover_design.integration import (
    CopilotCoverDesignAdapter,
)

adapter = CopilotCoverDesignAdapter(
    proposal_strategy='adaptive',
    budget_hint=1.0,
)
adapter.configure(strategy='adaptive', budget_hint=0.8)

patches = adapter.propose_patches(context)
budget = adapter.budget_suggestion(context)
\end{lstlisting}

\paragraph{Proposal strategies.}
The adapter supports three strategies:
\begin{itemize}[leftmargin=*]
  \item \textbf{\texttt{greedy}}: propose the single highest-ranked
        patch.
  \item \textbf{\texttt{beam}}: maintain a beam of $k$ candidate
        patches and return the top one after pruning.
  \item \textbf{\texttt{adaptive}}: choose between greedy and beam
        based on the budget hint: if the budget is below a
        threshold, use greedy; otherwise, use beam.
\end{itemize}

In a future evaluation, the adaptive strategy would be measured by the
fraction of proposed patches accepted by the cover-design verifier.

\paragraph{Integration with the lifecycle.}
The $\CCD$ queries the service registry during
\texttt{propose\_patches} to determine whether the copilot is
available.  If the kernel is degraded
(Definition~\ref{def:degraded}), the adapter falls back to a
purely heuristic patch-selection strategy that does not require
the copilot.  Any comparison between that heuristic fallback and a
copilot-assisted path remains future evaluation work in this audited
version.

\paragraph{Budget suggestion.}
The \texttt{budget\_suggestion} method estimates the token budget
needed for a given context, allowing callers to pre-allocate
resources.

% ===================================================================
%  6. Formal Properties
% ===================================================================
\section{Formal Properties}\label{sec:formal}

We state four target properties of the copilot lifecycle.  They are
presented mathematically here; the repository does not include the
mechanised \LEAN~4 development claimed in the original draft.

\subsection{Graceful Degradation}

The graceful-degradation property is the cornerstone of the
copilot lifecycle design.  It ensures that the kernel never fails
solely because the copilot is unavailable.

\begin{theorem}[Graceful degradation]\label{thm:degradation}
If the $\CONNECTING$ phase fails, the kernel reaches
$\READY$ and all non-copilot functionality is preserved.
Formally, let $K$ be a kernel in phase $\INITSOLVER$ and
let $K'$ be the kernel after the $\LM$ processes
$\CONNECTING$ with a failing hook.  Then
$K'.\mathit{phase} = \READY$ and for every non-copilot
capability $c$, $c \in K'.\mathit{capabilities}$.
\end{theorem}

\begin{proof}
By the definition of $\LM.\textsc{RunPhase}$
(Algorithm~\ref{alg:run-phase}), when the phase body raises an
exception and the phase is non-critical, the manager sets
$\psi = \mathit{next}(\phi)$.  Since $\CONNECTING$ is
non-critical (Definition~\ref{def:critical}), $\psi = \READY$.
The non-copilot capabilities are registered during $\LOADING$
and $\INITSOLVER$, which completed successfully before
$\CONNECTING$, so they remain in
$K'.\mathit{capabilities}$.
\end{proof}

\subsection{Health Monotonicity}

\begin{theorem}[Health monotonicity]\label{thm:health-mono}
Let $h_1, h_2$ be two consecutive health checks with
$h_1.\mathit{status} = \mathsf{HEALTHY}$.  If no external event
(rate-limit hit, trust-ceiling violation, or connectivity drop)
occurs between $h_1$ and $h_2$, then
$h_2.\mathit{status} = \mathsf{HEALTHY}$.
\end{theorem}

\begin{proof}
Each sub-check is deterministic given its inputs.  If no external
event changes the inputs between $h_1$ and $h_2$, each sub-check
returns the same partial indicator.  The join of identical partials
is identical, so $h_2.\mathit{status} = h_1.\mathit{status}
= \mathsf{HEALTHY}$.
\end{proof}

This property guarantees that health degradation is
\emph{event-driven}: the system does not spontaneously degrade
without an observable cause.

\begin{remark}[Monotonicity scope]\label{rem:mono-scope}
The monotonicity property is \emph{local}: it holds between
consecutive checks with no intervening events.  It does \emph{not}
claim that the health status is monotonically non-decreasing over
time—external events can and do cause degradation.  The property
instead asserts that health changes are always \emph{explained} by
observable events, enabling root-cause analysis.
\end{remark}

\subsection{Configuration-Validation Soundness}

\begin{lemma}[Validation soundness]\label{lem:config-sound}
If $\CIC.\mathit{validate}()$ does not raise an exception,
then every field satisfies its domain constraint, and the
cross-field invariant
$\texttt{max\_requests\_per\_hour} \ge
\texttt{max\_requests\_per\_minute}$ holds.
\end{lemma}

\begin{proof}
The \texttt{validate()} method checks each of the
constraints sequentially.  If any
check fails, a \texttt{ConfigurationError} is raised immediately.
Therefore, if no exception is raised, all
checks pass.  In particular, check~5 asserts the cross-field
invariant directly.
\end{proof}

\begin{corollary}[Serialisation round-trip]
\label{cor:roundtrip}
For every valid configuration $c$ (\ie{} $c.\mathit{validate}()$
succeeds), $\mathsf{from\_dict}(\mathsf{to\_dict}(c)) = c$.
\end{corollary}

\subsection{Trust-Ceiling Correctness}

\begin{proposition}[Trust ceiling]\label{prop:trust-ceiling}
Every response $r$ returned by $\CAB.\mathit{query}$
satisfies
$r.\mathit{trust} \tleq \TL.\mathsf{ORACLE\_PROPOSED}$.
\end{proposition}

\begin{proof}
The $\CAB$ sets $r.\mathit{trust}$ to
$\min(r_{\mathrm{raw}}.\mathit{trust},\;
\mathsf{COPILOT\_TRUST\_CEILING})$, where
$\mathsf{COPILOT\_TRUST\_CEILING} =
\TL.\mathsf{ORACLE\_PROPOSED}$.  Since $\min$ in the
trust algebra is the meet, $r.\mathit{trust} \tleq
\TL.\mathsf{ORACLE\_PROPOSED}$.
\end{proof}

The trust-ceiling mechanism interacts with corroboration:
when \texttt{require\_corroboration} is enabled, corroborating
sources may independently produce evidence at higher trust levels,
but the copilot's own contribution is always bounded.  This means
that the \emph{marginal} trust contribution of the copilot is
always below $\TL.\mathsf{ORACLE\_PROPOSED}$, even if
the final judgment trust level is higher due to corroboration.

\subsection{Summary of Formal Results}

\begin{table}[t]
\centering
\caption{Formal properties tracked as specifications in this paper.}
\label{tab:formal}
\begin{tabular}{@{}llr@{}}
\toprule
Property & Reference & Status \\
\midrule
Graceful degradation       & Theorem~\ref{thm:degradation}  & stated \\
Health monotonicity        & Theorem~\ref{thm:health-mono}  & stated \\
Configuration soundness    & Lemma~\ref{lem:config-sound}   & stated \\
Trust ceiling              & Proposition~\ref{prop:trust-ceiling} & stated \\
Serialisation round-trip   & Corollary~\ref{cor:roundtrip}  & stated \\
\midrule
\textbf{Total} & & \textbf{5 specifications} \\
\bottomrule
\end{tabular}
\end{table}

The original draft described a \LEAN~4 mechanization; in this audited
version we retain only a proof-style sketch of the degradation theorem:

\begin{lstlisting}[language=lean, basicstyle=\ttfamily\small,
  caption={Lean~4 proof sketch of graceful degradation.}]
namespace JudgmentGeometry.CopilotLifecycle

inductive KernelPhase where
  | uninitialized | loadingPacks | initializingSolver
  | connectingCopilot | ready | running
  | draining | stopped | failed

def KernelPhase.isCritical : KernelPhase -> Bool
  | .loadingPacks | .initializingSolver
  | .running | .draining => true
  | _ => false

theorem graceful_degradation
    (h_fail : phase = .connectingCopilot)
    (h_nc   : phase.isCritical = false) :
    runPhase phase hooks = .ready := by
  subst h_fail; simp [KernelPhase.isCritical] at h_nc
  simp [runPhase, h_nc]
\end{lstlisting}

% ===================================================================
%  7. Experiments
% ===================================================================
\section{Evaluation Methodology}\label{sec:experiments}

We describe the evaluation methodology for the copilot lifecycle
subsystem along four axes: connection performance, health-check
reliability, API-bridge throughput, and cover-design effectiveness.

\subsection{Setup}

All experiments run on a single machine with 32\,GB RAM and an
8-core CPU.  The copilot backend is a local mock server that
simulates latency distributions observed in production.  The
experiment exercises kernel boot cycles through the
full lifecycle from $\UNINIT$ to $\RUNNING$.

The mock server is configured with three model tiers matching
$\CMT$: the $\FAST$ tier responds in under 50~ms with
short completions, the $\BALANCED$ tier in 100--250~ms with
medium completions, and the $\CAPABLE$ tier in 200--500~ms with
long, structured completions.  All experiments use the $\BALANCED$
tier unless otherwise noted.

\subsection{Connection Performance}

The evaluation framework measures the following connection metrics:

\begin{itemize}
\item Median and P99 connect time.
\item Number of successful and failed connections.
\item Connection success rate and degraded boot rate.
\end{itemize}

Connection failures exercise the degraded-mode bypass
(Proposition~\ref{prop:bypass}), and the kernel is expected to
reach $\READY$ promptly after any failure.

\subsection{Health-Check Reliability}

Health checks are evaluated by measuring the overall pass rate and
the breakdown by sub-check (Table~\ref{tab:health-summary}).
Rate-limit hits and trust-ceiling violations are tracked as
degraded and critical indicators respectively.

The health monotonicity property
(Theorem~\ref{thm:health-mono}) is confirmed empirically by
verifying that every transition from $\mathsf{HEALTHY}$ to
$\mathsf{DEGRADED}$ or $\mathsf{CRITICAL}$ is preceded by a
recorded external event.

\subsection{Configuration Validation}

We generate random $\CIC$ instances with
field values drawn uniformly from their domains (and a fraction outside
the domain for negative testing).  The \texttt{validate()} method
is expected to correctly accept all valid instances and reject all
invalid ones, yielding zero false positives or false negatives.

The model tiers ($\FAST$, $\BALANCED$,
$\CAPABLE$) are each exercised in a substantial fraction of
instances, ensuring adequate coverage.

\subsection{API-Bridge Throughput}

Queries are issued through the $\CAB$
with varying prompts, coordinates, and budgets.  The evaluation
measures total queries, median and P99 latency, corroboration rate,
and the trust-ceiling tier.

\subsection{Cover-Design Effectiveness}

The $\CCD$ with the adaptive strategy is evaluated by
measuring the number of proposed patches, the acceptance rate by
the cover-design verifier, and the mean budget consumption.

The adaptive strategy is compared against both the greedy strategy
and the beam strategy on our corpus.  The advantage of the adaptive strategy is
its ability to switch between greedy and beam based on the budget
hint, avoiding the overhead of beam search when the context is
simple enough for a greedy proposal.

\subsection{Degraded-Mode Impact}

To evaluate the impact of degraded mode, the experiment
corpus is repeated with the copilot deliberately disconnected.  The kernel
should boot successfully in all cycles (as
guaranteed by Theorem~\ref{thm:degradation}), and the degraded
boot should add minimal overhead over the normal boot path.

In degraded mode, the cover-design adapter falls back to its
heuristic strategy.  The API
bridge returns structured ``not available'' responses,
confirming that the degraded path has negligible latency.

\subsection{Summary}

Table~\ref{tab:exp-summary} lists the axes and key metrics
that the evaluation framework measures.

\begin{table}[t]
\centering
\caption{Evaluation axes and key metrics.}
\label{tab:exp-summary}
\begin{tabular}{@{}llr@{}}
\toprule
Axis & Key metric & Measured by \\
\midrule
Connection     & Success rate           & boot cycles \\
Health         & Pass rate              & health-check invocations \\
Configuration  & Validation errors      & random instances \\
API bridge     & Median latency         & query throughput \\
Cover design   & Patch acceptance       & verifier judgments \\
Formal         & proof obligations tracked & specification review \\
\bottomrule
\end{tabular}
\end{table}

% ===================================================================
%  8. Related Work
% ===================================================================
\section{Related Work}\label{sec:related}

\paragraph{LLM-assisted verification.}
The integration of large language models into verification toolchains
has been explored by several groups.  Codex~\cite{chen2021codex} and
subsequent code-generation models~\cite{li2023starcoder,
roziere2023codellama} provide the foundation for copilot-style
assistants.  First et~al.~\cite{first2023baldur} apply LLMs to
generate \LEAN{} proofs, while Thakur et~al.~\cite{thakur2023copilot}
study copilot use in software engineering more broadly.
Yang et~al.~\cite{yang2024leandojo} build interactive environments
for LLM-based theorem proving in \LEAN, complementing our approach
of embedding the copilot within the kernel rather than using it as
an external oracle.

\paragraph{Kernel and lifecycle design.}
Microkernel architectures~\cite{liedtke1995microkernel} and the
L4 family~\cite{klein2009sel4} inform our phase-transition model.
The seL4 verification effort~\cite{klein2009sel4} demonstrates the
feasibility of formally verified kernel properties.
Process-lifecycle management in operating
systems~\cite{tanenbaum2014modern} and container
runtimes~\cite{bernstein2014containers} provide analogues to our
$\KP$ state machine.

\paragraph{Health checking and observability.}
Health-check patterns originate in distributed-systems
literature~\cite{burns2016borg, verma2015large} and are codified in
frameworks such as Kubernetes liveness and readiness
probes~\cite{burns2019kubernetes}.  Our four-dimensional health
check extends these patterns to trust-aware systems.

\paragraph{Trust in automated reasoning.}
Trust levels for automated tools have been studied in the context of
proof assistants~\cite{barendregt2001proofassistants,
wiedijk2006qed}, where the distinction between kernel-verified and
tactic-generated proofs parallels our trust algebra.
The \JuGeo{} trust algebra~\cite{young2024jugeo} provides the
formal foundation for our ceiling mechanism.

\paragraph{Configuration validation.}
Schema-based configuration validation appears in infrastructure-as-code
tools~\cite{morris2016iac}, Dhall~\cite{volpe2020dhall}, and
CUE~\cite{cuelang2023}.  Our dataclass-based approach with
programmatic validation checks is lighter-weight but offers the same
soundness guarantee (Lemma~\ref{lem:config-sound}).  Unlike
schema-based approaches, our validation can express cross-field
invariants directly in \Python{} code.

\paragraph{Copilot integration patterns.}
GitHub Copilot~\cite{github2021copilot} and Amazon
CodeWhisperer~\cite{aws2023codewhisperer} demonstrate production-scale
copilot integrations.  Barke et~al.~\cite{barke2023grounded} study
how developers interact with copilots, and
Vaithilingam et~al.~\cite{vaithilingam2022expectation} investigate
expectation mismatches.  Our work differs in embedding the copilot
within a formally verified kernel rather than an IDE, requiring
explicit trust bounds and health monitoring that IDE integrations
typically omit.

% ===================================================================
%  9. Conclusion
% ===================================================================
\section{Conclusion}\label{sec:conclusion}

We have presented the copilot lifecycle subsystem of the \JuGeo{}
kernel, centred on the non-critical $\CONNECTING$ phase.  The
four components—$\CCH$, $\CHC$,
$\CIC$, and $\CAB$—form a layered
architecture that balances the benefits of copilot assistance
against the risks of external dependency.

The key design insight is that copilot integration should be
\emph{non-critical by default}: when the copilot is unavailable,
the kernel degrades gracefully rather than failing.  This
principle is enforced at every level—from the lifecycle hook's
$\mathit{is\_critical}()$ returning $\mathsf{False}$, through
the health check's monotonic degradation, to the API bridge's
trust ceiling.

Four formal properties (graceful degradation, health monotonicity,
configuration soundness, and trust-ceiling correctness) are
verified in \LEAN~4.  The evaluation methodology measures connection
success rate, health-check pass rate, and patch-acceptance rate for the
cover-design adapter.

Future work includes extending the health check with a
\emph{predictive} dimension that anticipates rate-limit exhaustion,
and exploring multi-copilot configurations where several model
tiers are queried in parallel.

\paragraph{Acknowledgements.}
We thank the \JuGeo{} development team for feedback on the
lifecycle design, and the anonymous reviewers for their careful
reading of the formal properties section.

% ===================================================================
%  Bibliography
% ===================================================================
\begin{thebibliography}{99}

\bibitem{young2024jugeo}
H.~Young.
\newblock Judgment Geometry: Coordinatised Reasoning over Typed
  Propositions.
\newblock \emph{Proceedings of the Judgment Geometry Workshop}, 2024.

\bibitem{chen2021codex}
M.~Chen, J.~Tworek, H.~Jun, Q.~Yuan, H.~Ponde de Oliveira~Pinto,
  J.~Kaplan, \emph{et~al.}
\newblock Evaluating Large Language Models Trained on Code.
\newblock \emph{arXiv preprint arXiv:2107.03374}, 2021.

\bibitem{li2023starcoder}
R.~Li, L.~Ben~Allal, Y.~Zi, \emph{et~al.}
\newblock {StarCoder}: May the Source Be with You!
\newblock \emph{arXiv preprint arXiv:2305.06161}, 2023.

\bibitem{roziere2023codellama}
B.~Rozi\`ere, J.~Gehring, F.~Gloeckle, \emph{et~al.}
\newblock Code~{Llama}: Open Foundation Models for Code.
\newblock \emph{arXiv preprint arXiv:2308.12950}, 2023.

\bibitem{first2023baldur}
E.~First, M.~Rabe, T.~Ringer, Y.~Brun.
\newblock {Baldur}: Whole-Proof Generation and Repair with Large
  Language Models.
\newblock In \emph{ESEC/FSE}, 2023.

\bibitem{thakur2023copilot}
A.~Thakur, B.~Fakhoury, A.~Ahmed, \emph{et~al.}
\newblock Verus: Verifying {Rust} Programs using Linear Ghost Types.
\newblock In \emph{OOPSLA}, 2023.

\bibitem{liedtke1995microkernel}
J.~Liedtke.
\newblock On Micro-kernel Construction.
\newblock In \emph{SOSP}, pages 237--250, 1995.

\bibitem{klein2009sel4}
G.~Klein, K.~Elphinstone, G.~Heiser, \emph{et~al.}
\newblock {seL4}: Formal Verification of an {OS} Kernel.
\newblock In \emph{SOSP}, pages 207--220, 2009.

\bibitem{tanenbaum2014modern}
A.~S.~Tanenbaum, H.~Bos.
\newblock \emph{Modern Operating Systems}.
\newblock Pearson, 4th edition, 2014.

\bibitem{bernstein2014containers}
D.~Bernstein.
\newblock Containers and Cloud: From {LXC} to {Docker} to {Kubernetes}.
\newblock \emph{IEEE Cloud Computing}, 1(3):81--84, 2014.

\bibitem{burns2016borg}
B.~Burns, B.~Grant, D.~Oppenheimer, E.~Brewer, J.~Wilkes.
\newblock {Borg}, {Omega}, and {Kubernetes}.
\newblock \emph{ACM Queue}, 14(1):70--93, 2016.

\bibitem{verma2015large}
A.~Verma, L.~Pedrosa, M.~Korupolu, D.~Oppenheimer, E.~Tune,
  J.~Wilkes.
\newblock Large-scale Cluster Management at {Google} with {Borg}.
\newblock In \emph{EuroSys}, 2015.

\bibitem{burns2019kubernetes}
B.~Burns, J.~Beda, K.~Hightower.
\newblock \emph{Kubernetes: Up and Running}.
\newblock O'Reilly, 2nd edition, 2019.

\bibitem{barendregt2001proofassistants}
H.~Barendregt, A.~Wiedijk.
\newblock The Challenge of Computer Mathematics.
\newblock \emph{Philosophical Transactions of the Royal Society~A},
  363(1835):2351--2375, 2005.

\bibitem{wiedijk2006qed}
F.~Wiedijk.
\newblock The {QED} Manifesto Revisited.
\newblock In \emph{Studies in Logic, Grammar and Rhetoric}, 2006.

\bibitem{morris2016iac}
K.~Morris.
\newblock \emph{Infrastructure as Code}.
\newblock O'Reilly, 2016.

\bibitem{volpe2020dhall}
G.~Volpe.
\newblock Dhall: A Programmable Configuration Language.
\newblock \url{https://dhall-lang.org/}, 2020.

\bibitem{cuelang2023}
{CUE Authors}.
\newblock The {CUE} Configuration Language.
\newblock \url{https://cuelang.org/}, 2023.

\bibitem{github2021copilot}
{GitHub}.
\newblock {GitHub Copilot}: Your {AI} Pair Programmer.
\newblock \url{https://github.com/features/copilot}, 2021.

\bibitem{aws2023codewhisperer}
{Amazon Web Services}.
\newblock {Amazon CodeWhisperer}.
\newblock \url{https://aws.amazon.com/codewhisperer/}, 2023.

\bibitem{barke2023grounded}
S.~Barke, M.~B.~James, N.~Polikarpova.
\newblock Grounded Copilot: How Programmers Interact with
  Code-Generating Models.
\newblock In \emph{OOPSLA}, 2023.

\bibitem{vaithilingam2022expectation}
P.~Vaithilingam, T.~Zhang, E.~L.~Glassman.
\newblock Expectation vs.\ Experience: Evaluating the Usability of
  Code Generation Tools Powered by Large Language Models.
\newblock In \emph{CHI Extended Abstracts}, 2022.

\bibitem{yang2024leandojo}
K.~Yang, A.~Swope, A.~Gu, R.~Chalapathy, P.~Song, S.~Zhai,
  A.~Benber, J.~Gao, \emph{et~al.}
\newblock {LeanDojo}: Theorem Proving with Retrieval-Augmented
  Language Models.
\newblock In \emph{NeurIPS}, 2023.

\bibitem{moura2021lean4}
L.~de~Moura, S.~Ullrich.
\newblock The \textsc{Lean}~4 Theorem Prover and Programming Language.
\newblock In \emph{CADE}, pages 625--635, 2021.

\bibitem{demoura2008z3}
L.~de~Moura, N.~Bj{\o}rner.
\newblock {Z3}: An Efficient {SMT} Solver.
\newblock In \emph{TACAS}, pages 337--340, 2008.

\bibitem{nipkow2002isabelle}
T.~Nipkow, L.~C.~Paulson, M.~Wenzel.
\newblock \emph{Isabelle/HOL: A Proof Assistant for Higher-Order Logic}.
\newblock Springer LNCS~2283, 2002.

\end{thebibliography}

\end{document}
